\chapter{The Pipeline: Data, Features, Models, Validation, Calibration} \label{chapter3}

This chapter describes the pipeline as it executes on the canonical run analysed in Chapter~\ref{chapter4}. Every numerical setting cited corresponds to a value stored in the system configuration; every structured output cited is produced as a side-effect of the run. The chapter is organised top-down: pipeline overview, universe and acquisition, feature engineering with the leakage audit, the five-model ensemble, the meta-learner, the calibration step, the walk-forward protocol, the backtest engine, and the parallelisation and reproducibility infrastructure.

\section{Pipeline Overview}

\begin{figure}[H]
\centering
\begin{adjustbox}{max width=\textwidth}
\begin{tikzpicture}[
    node distance=4mm and 7mm,
    every node/.style={font=\scriptsize\sffamily, rounded corners=2pt, align=center},
    stage/.style ={draw=pipelineblue,  fill=boxfill, thick, minimum width=34mm, minimum height=9mm},
    data/.style  ={draw=pipelinegreen, fill=boxfill, thick, minimum width=34mm, minimum height=9mm},
    model/.style ={draw=pipelineamber, fill=boxfill, thick, minimum width=34mm, minimum height=9mm},
    deploy/.style={draw=pipelinered,   fill=boxfill, thick, minimum width=34mm, minimum height=9mm},
    arrow/.style ={-Stealth, thick, draw=darkgray}
]
% Row 1: Acquisition (3 boxes side-by-side that converge)
\node[stage] (yf)   {yfinance OHLCV};
\node[stage, right=of yf]   (cues) {Global cues\\\tiny SP500, Nasdaq, VIX, DXY, Crude, Nikkei};
\node[stage, right=of cues] (cal)  {NSE calendar\\\tiny Expiry, RBI, Budget, Results};

% Row 2: raw store
\node[data,  below=of cues] (raw) {Raw OHLCV Store\\\tiny incremental per-ticker parquet};

% Row 3: feature pipeline
\node[stage, below=of raw] (feat) {Feature engine ($\sim$200 features)\\\tiny + leakage audit};
\node[stage, below=of feat] (scale){RobustScaler + PCA \emph{fit on train only}};

% Row 4: walk-forward
\node[stage, below=of scale] (wf) {Walk-forward, 6 folds, train 0.70$\to$0.95};

% Row 5: committee
\node[model, below=of wf, minimum width=120mm, minimum height=14mm] (committee)
  {\textbf{Five-model ensemble per (ticker, fold)}\\
   \tiny LightGBM\,$\cdot$\,XGBoost\,$\cdot$\,BiLSTM\,$\cdot$\,TCN-Transformer\,$\cdot$\,N-BEATS};

% Row 6: meta + temp
\node[stage, below=of committee] (meta) {Logistic meta-stack $\to$ temperature scaling $\sigma(z/T^\star)$};

% Row 7: gate
\node[stage, below=of meta] (gate) {Gate: $\hat p \geq 0.58$ \;AND\; $|\widehat\Delta| \geq 0.4\%$};

% Row 8: deployment chain (4 boxes, will fit since each is 34mm and gap 7mm => 4*34+3*7 = 157mm > 168mm? Adjust spacing)
\node[deploy, below=8mm of gate] (port)  {Portfolio backtest\\\tiny top-15 cross-sectional};
\node[deploy, right=of port]     (paper) {Paper trader\\\tiny Indian cost model};
\node[deploy, below=of port]     (promo) {Promotion gate\\\tiny Sharpe, DD, Brier};
\node[deploy, right=of promo]    (live)  {Angel One SmartAPI\\\tiny live execution};

% Arrows
\draw[arrow] (yf.south)   |- (raw.west);
\draw[arrow] (cues) -- (raw);
\draw[arrow] (cal.south)  |- (raw.east);
\draw[arrow] (raw)   -- (feat);
\draw[arrow] (feat)  -- (scale);
\draw[arrow] (scale) -- (wf);
\draw[arrow] (wf)    -- (committee);
\draw[arrow] (committee) -- (meta);
\draw[arrow] (meta)  -- (gate);
\draw[arrow] (gate)  -- (port);
\draw[arrow] (port)  -- (paper);
\draw[arrow] (paper.south) -- (live.north);
\draw[arrow] (port)  -- (promo);
\draw[arrow] (promo) -- (live);

\end{tikzpicture}
\end{adjustbox}
\caption{End-to-end pipeline overview.}
\label{fig:pipeline_overview}
\end{figure}

\section{Universe and Acquisition}

\subsection{Universe}

The configured universe is the Nifty-100 (Nifty-50 + Nifty Next-50) of the National Stock Exchange of India, grouped into 14 sector buckets (banking, IT, automobiles, FMCG, pharma, energy, metals, telecom, capital goods, cement, consumer durables, real estate, conglomerate / infra, defence). On the latest full run (May 2026), all $100$ tickers passed the post-feature minimum-history filter. Table~\ref{tab:universe_sectors} reports the sector distribution.

\begin{table}[H]
\centering
\caption{Nifty-100 sector distribution}
\label{tab:universe_sectors}
\renewcommand{\arraystretch}{1.3}
\begin{adjustbox}{max width=\textwidth}
\begin{tabular}{l c p{9cm}}
\toprule
\textbf{Sector} & \textbf{Count} & \textbf{Representatives} \\
\midrule
Banking \& Finance   & 21 & SBIN, HDFCBANK, ICICIBANK, AXISBANK, KOTAKBANK, INDUSINDBK, BAJFINANCE, BAJAJFINSV, HDFCLIFE, SBILIFE, ICICIGI, MUTHOOTFIN, CHOLAFIN, SHRIRAMFIN, BANDHANBNK, IDFCFIRSTB, FEDERALBNK, AUBANK, RBLBANK, MANAPPURAM, LICHSGFIN \\
IT                   & 12 & TCS, INFY, WIPRO, HCLTECH, TECHM, LTIM, MPHASIS, PERSISTENT, COFORGE, TATAELXSI, OFSS, NAUKRI \\
Automobiles          & 8  & MARUTI, M\&M, TVSMOTOR, BAJAJ-AUTO, HEROMOTOCO, EICHERMOT, MOTHERSON, EXIDEIND \\
FMCG                 & 8  & HINDUNILVR, ITC, NESTLEIND, BRITANNIA, TATACONSUM, MARICO, COLPAL, GODREJCP \\
Pharma               & 8  & SUNPHARMA, DRREDDY, CIPLA, DIVISLAB, LUPIN, TORNTPHARM, AUROPHARMA, ALKEM \\
Energy               & 8  & RELIANCE, ONGC, BPCL, NTPC, POWERGRID, COALINDIA, GAIL, TATAPOWER \\
Metals \& Mining     & 6  & TATASTEEL, HINDALCO, JSWSTEEL, VEDL, SAIL, NMDC \\
Telecom              & 2  & BHARTIARTL, INDUSTOWER \\
Capital Goods        & 8  & LT, BHEL, SIEMENS, ABB, HAVELLS, POLYCAB, CUMMINSIND, BOSCHLTD \\
Cement               & 4  & ULTRACEMCO, GRASIM, AMBUJACEM, SHREECEM \\
Consumer Durables    & 7  & TITAN, ASIANPAINT, PIDILITIND, BERGEPAINT, VOLTAS, PAGEIND, DMART \\
Real Estate          & 2  & DLF, GODREJPROP \\
Conglomerate / Infra & 4  & ADANIENT, ADANIPORTS, IRFC, ETERNAL \\
Defence              & 2  & BEL, HAL \\
\midrule
\textbf{Total}       & \textbf{100} & 99 trained, 1 skipped (BAJAJHFL) \\
\bottomrule
\end{tabular}
\end{adjustbox}
\end{table}

\subsection{Acquisition}

Per-ticker OHLCV is downloaded from Yahoo Finance~\cite{yfinance} with a 0.3-second inter-request delay, parallelised through an 8-worker thread pool (I/O-bound). Each ticker's history is stored as an incremental parquet: subsequent runs read the existing record, identify the latest date, and fetch only new rows. Two global datasets are refreshed once per run: (i) global market cues comprising S\&P~500, Nasdaq, US VIX, DXY, Crude, Nikkei, Nifty-50, and Nifty Bank; and (ii) the USD/INR exchange rate. Data start date is 1 January 2018.

\section{Feature Engineering}

\subsection{Feature Families}

The feature engine computes \(\sim\)200 features per row from the OHLCV plus the merged exogenous cues. The families are summarised in Table~\ref{tab:feature_families}. The post-engineering matrix is reduced by per-fold feature selection to 50 features; on the canonical run, the per-ticker selected count is $47$ or $50$ depending on whether the ticker belongs to the IT sector (where USD/INR and Nasdaq features are force-included) or the general sector bucket.

\begin{table}[H]
\centering
\caption{Feature families in the pipeline}
\label{tab:feature_families}
\renewcommand{\arraystretch}{1.4}
\begin{adjustbox}{max width=\textwidth}
\begin{tabularx}{\textwidth}{l X}
\toprule
\textbf{Family} & \textbf{Members (representative)} \\
\midrule
Returns and ratios       & 1/3/5/10/20-day arithmetic and log returns; OHLC ratios; gap features; cumulative drawdown; rolling sum-of-squares \\
Trend                    & SMA/EMA at $\{5,10,20,50,100,200\}$, distance to MA, MA crossovers, MACD (12/26/9), MACD histogram \\
Mean-reversion           & RSI(14), RSI(7), z-score of return at $\{5,10,20\}$, distance to 52-week high/low, Bollinger band z-score \\
Volatility               & ATR(14), realised vol at $\{5,10,20\}$, Parkinson, Garman--Klass, Yang--Zhang estimators, volatility z-score \\
Volume                   & OBV, CMF(20), MFI(14), volume moving averages, volume z-score, dollar-volume \\
Range / candlestick      & Body / shadow ratios, doji / hammer / engulfing flags, true range, mid-day pivot \\
Calendar                 & Day-of-week, month, quarter (cyclic), \texttt{days\_to\_expiry}, \texttt{is\_expiry\_week}, \texttt{is\_expiry\_day}, \texttt{days\_to\_rbi}, \texttt{is\_rbi\_week}, \texttt{days\_to\_budget}, \texttt{is\_budget\_week}, \texttt{is\_result\_season} \\
Global cues (shift +1d)  & \texttt{sp500\_ret\_prev/5d}, \texttt{nasdaq\_ret\_prev/5d}, \texttt{us\_vix\_level/zscore/spike}, \texttt{dxy\_ret\_prev/5d}, \texttt{crude\_ret\_prev/5d}, \texttt{nikkei\_ret\_prev}, \texttt{nifty50\_ret\_prev/5d}, \texttt{niftybank\_ret\_prev/5d/20d} \\
Sector-specific          & IT: USD/INR returns, RSI, 20-day MA ratio, alpha vs.\ INR. Banking: days to next RBI meeting, Nifty Bank momentum. Force-included per sector. \\
Regime tags              & 3-state regime label (bear / sideways / bull) derived from rolling return percentiles \\
\bottomrule
\end{tabularx}
\end{adjustbox}
\end{table}

\subsection{The Target}

The classification target is binary up/down at horizon $h = 1$, but with a dead-band:
\[
    y_t \;=\; \mathbb{1}\!\left[\, \frac{C_{t+1} - C_t}{C_t} \;\geq\; \delta \,\right], \qquad \delta = 0.004.
\]
Days with $|r_{t+1}| < 0.004$ are retained in training and labelled by their sign; the minimum-move threshold $\delta$ is applied at signal time, not at label time. The shift is performed once, \emph{after} all features have been computed and before any train/test split.

\subsection{Leakage Audit}

Table~\ref{tab:leakage_audit} documents the five places where look-ahead could enter and how the pipeline prevents it.

\begin{table}[H]
\centering
\caption{Leakage audit and corresponding pipeline controls}
\label{tab:leakage_audit}
\renewcommand{\arraystretch}{1.25}
\footnotesize
\begin{tabularx}{\textwidth}{@{}>{\raggedright\arraybackslash}p{2.6cm} >{\raggedright\arraybackslash}X >{\raggedright\arraybackslash}X@{}}
\toprule
\textbf{Risk} & \textbf{Where it would enter} & \textbf{Pipeline control} \\
\midrule
Target shift              & A feature shifted by $-N$ would share a future row with the target. & Target built by a single forward label shift applied after all features. No feature uses a forward shift. \\
Scaler fit                & Fitting a scaler on the whole panel leaks the test-period mean into train. & Robust scaler fit on the training window only; transform applied to validation and test. \\
PCA fit                   & Same as scaler.                                                       & PCA fit on the training window only. \\
Global cues               & US close of date $t$ is published during the Indian session of $t$. & US-side cues shifted by one calendar day, then merged by backward temporal join to the Indian daily series. \\
Feature selection         & Selecting on the full panel by importance on the full target is a peek. & Top-50 features chosen per fold on the training slice only; sector-specific force-include lists do not see the target. \\
Temperature scaling       & Fitting $T$ on the test set would inflate the gate.                  & $T^\star$ fit on the validation slice only; applied unchanged at inference. \\
\bottomrule
\end{tabularx}
\end{table}

\section{FinBERT-India: Domain-Adapted Sentiment Module}
\label{sec:finbert_india}

A common gap in the Indian-market literature is the absence of a sentiment signal channel: prior studies treat the problem as purely OHLCV-driven. To address this gap, a domain-adapted sentiment encoder, FinBERT-India, was constructed; the trained checkpoint ships with the accompanying source tree. The module is described in this section for completeness, but it is important to state up front that the sentiment channel has \textbf{not} yet been integrated into the walk-forward feature pipeline, and therefore did \emph{not} contribute to the May 2026 canonical results reported in Chapter~\ref{chapter4}. The system as evaluated is a pure price, macro, and calendar baseline; the sentiment integration is documented as immediate future work in Chapter~\ref{chapter6}.

\subsection{Base Model and Domain Adaptation Rationale}

The base model is \texttt{ProsusAI/finbert}~\cite{araci2019finbert}, a BERT-base encoder pre-trained on English financial corpora (Financial PhraseBank, Reuters TRC2-financial). FinBERT outperforms general-purpose BERT on English financial sentiment but, out of the box, is mis-calibrated on Indian-market text: it does not handle India-specific vocabulary such as \enquote{NPA}, \enquote{CASA ratio}, \enquote{FII outflow}, \enquote{RBI MPC verdict}, \enquote{rupee depreciation}, or Indian-broker phrasings of earnings beats and misses. The objective of fine-tuning is to transfer the financial-sentiment prior of FinBERT to the Indian-market lexicon while preserving the original three-way label scheme.

\subsection{Training Corpus}

A hand-curated corpus of \textbf{687 labelled Indian-finance sentences} was compiled, balanced across ten sectors: banking, IT, pharma, FMCG, auto, metals, telecom, power, real estate, and oil~\&~gas. Each sentence is a single financial-news headline-style sentence of 4--15 words (mean 10.7) labelled into one of three classes by domain-aware annotation:
\begin{itemize}\itemsep1pt
    \item \textbf{Positive (bullish)}: 198 samples (28.8\%). Examples: \enquote{HDFC Bank Q4 profit jumps 23 percent year on year beating estimates}; \enquote{TCS Q4 revenue beats estimates deal wins at record 11 billion dollars}.
    \item \textbf{Negative (bearish)}: 238 samples (34.6\%). Examples: \enquote{SBI gross NPA surges to 5.1 percent from 3.8 percent previous quarter}; \enquote{Infosys cuts FY guidance citing weak banking sector spending in US}.
    \item \textbf{Neutral (informational)}: 251 samples (36.5\%). Examples: \enquote{HDFC Bank board meeting scheduled April 18 to consider Q4 FY26 results}; \enquote{TCS to announce quarterly results on April 10 2026}.
\end{itemize}

The corpus deliberately concentrates on phrasings that are specific to Indian market reportage: result-day verdicts, RBI MPC outcomes, F\&O expiry commentary, FII/DII flow updates, asset-quality ratios, and rupee-INR moves. These are exactly the phrasings that an off-the-shelf English-trained FinBERT mis-classifies.

\subsection{Fine-Tuning Procedure}

Training was performed with the HuggingFace \texttt{Trainer} API. Key details:

\begin{itemize}\itemsep1pt
    \item \textbf{Split}: 80\%/20\% stratified by class, fixed seed 42; $n_{\text{train}}=549$, $n_{\text{val}}=138$.
    \item \textbf{Optimisation}: AdamW with learning rate $1\times10^{-5}$, weight decay $0.01$, batch size 16, 8 epochs.
    \item \textbf{Label remapping}: the training corpus uses the ordering $\{0\!=\!\text{neg},\,1\!=\!\text{neu},\,2\!=\!\text{pos}\}$ but FinBERT's pre-trained classification head outputs $\{0\!=\!\text{pos},\,1\!=\!\text{neg},\,2\!=\!\text{neu}\}$. Training labels are remapped to FinBERT's order before training so that the pre-trained classification head is preserved rather than randomly re-initialised.
    \item \textbf{Tokeniser}: the original FinBERT WordPiece tokeniser (max length 128 tokens) is retained without modification.
\end{itemize}

\subsection{Validation Performance}

On the held-out validation split (138 sentences), the fine-tuned FinBERT-India model attains:
\begin{itemize}\itemsep1pt
    \item Validation accuracy: \textbf{84.06\%}
    \item Macro F1-score: \textbf{84.10\%}
\end{itemize}

These metrics are recorded alongside the trained weights in the repository for reproducibility. The accuracy comfortably exceeds the random baseline ($\sim 33\%$ for three balanced classes), confirming that the model has learnt non-trivial Indian-finance sentiment structure from the small but domain-focused fine-tuning corpus.

\subsection{Inference Path and Backfill}

A companion backfill script consumes Yahoo Finance news headlines per ticker, runs each headline through FinBERT-India, and writes a daily-aggregated sentiment record per ticker to a sentiment-history store. Each daily record contains a positive--neutral--negative probability distribution and a single signed sentiment score per ticker; these are designed as drop-in features that, once integrated into the walk-forward feature pipeline (Section~\ref{sec:walkforward}), will be subject to the same per-fold scaler, PCA, and leakage-audit discipline as the OHLCV and macro features. The marginal contribution of the sentiment channel to OOS accuracy and to the gated bootstrap accuracy is the immediate next experiment.

\section{The Five-Model Ensemble}

The system trains five heterogeneous binary classifiers per (ticker, fold). The tree pair (LightGBM, XGBoost) consumes a two-dimensional PCA-projected matrix; the three deep-sequence models (BiLSTM, TCN-Transformer, N-BEATS) consume a three-dimensional tensor with sequence length $T=20$ built by a sliding window over the scaled (non-PCA) features. The ensemble specification is in Table~\ref{tab:model_committee}.

\begin{table}[H]
\centering
\caption{The five-model ensemble}
\label{tab:model_committee}
\renewcommand{\arraystretch}{1.4}
\begin{adjustbox}{max width=\textwidth}
\begin{tabularx}{\textwidth}{l l X}
\toprule
\textbf{Family} & \textbf{Model} & \textbf{Topology and key hyper-parameters} \\
\midrule
Trees   & LightGBM         & 1000 trees, depth 5, lr $0.01$, leaves 31, subsample 0.8, colsample 0.8, $\alpha=0.3$, $\lambda=1.5$, min child 20, early stopping 50 rounds. \\
Trees   & XGBoost          & 1000 trees, depth 5, lr $0.01$, subsample 0.8, colsample 0.8, $\alpha=0.3$, $\lambda=1.5$, early stopping 50 rounds, logistic objective. \\
RNN     & BiLSTM           & 2-layer bidirectional (32$\to$16 per direction, merged $\to$ 64$\to$32), dropout 0.3, recurrent dropout 0.2, $\ell_2=10^{-4}$, Adam lr $10^{-3}$, batch 32, max 100 epochs, ES patience 8. Safe under leakage audit as all input features are pre-lagged. \\
TCN+Attn & TCN-Transformer & TCN (dilated causal convolutions, dilations $\{1,2,4\}$, kernel 3, $d_\mathrm{model}=64$) followed by 4-head self-attention (key dim 16). Receptive field $\approx 29$ steps. \\
N-BEATS & N-BEATS          & 3 residual backcast blocks, 4 FC layers per block, $\mathrm{fc\_dim}=512$, projection $\mathrm{proj\_dim}=256$, sigmoid classification head. Adam lr $5 \times 10^{-4}$. \\
\bottomrule
\end{tabularx}
\end{adjustbox}
\end{table}

Storage budget is preserved by saving deep-learning model weights only for the \emph{last} fold per ticker (used for the production copy and live inference); intermediate-fold weights are discarded. This reduces checkpoint storage from $N_\mathrm{folds} \times N_\mathrm{DL}$ per ticker to $1 \times N_\mathrm{DL}$ per ticker, a $\sim$83\% reduction.

\section{Meta-Learner and Calibration}

\subsection{Logistic Stacking}

For each fold, the LightGBM and XGBoost out-of-fold validation probabilities are stacked into a 2-feature design matrix; a logistic regression ($\ell_2$ regularisation, $C=1$) is fit to map these probabilities to the validation labels. The three deep-learning probabilities are admitted into the stack only if their validation accuracy clears a sanity floor (otherwise their column is dropped to avoid feeding pure noise to the meta-learner). On the canonical run, the deep-learning pass-through floor was not consistently met across tickers, so the dominant meta-signal flowed through the tree stack; the resulting meta-probability therefore primarily reflects the logistic blend of LightGBM and XGBoost.

\subsection{Temperature Scaling}

Let $p_v$ be the meta-probability on the validation slice and $z_v = \mathrm{logit}(p_v)$. We fit
\[
T^\star \;=\; \arg\min_{T \in [0.05, 5]} \; \mathrm{NLL}\!\left(\sigma(z_v / T),\ y_v\right),
\]
using one-dimensional scalar optimisation, and apply $\hat p = \sigma(\mathrm{logit}(p) / T^\star)$ to the test slice. Temperature scaling preserves the model's ranking (so the AUC and the argmax accuracy are unchanged); it shifts the entire probability distribution toward observed frequencies so that the gate threshold $0.58$ becomes operationally meaningful.

\section{Walk-Forward Validation}\label{sec:walkforward}

\subsection{Protocol}

For each ticker with $N$ rows after feature alignment, the walk-forward protocol generates six chronological folds. Fold $k$ ($k=1,\dots,6$) uses
\[
    \mathrm{train\_ratio}_k \;=\; 0.70 + 0.05\,(k-1), \quad
    \mathrm{train\_size}_k \;=\; \lfloor N \cdot \mathrm{train\_ratio}_k \rfloor,
\]
with the next $10\%$ of rows used as validation and the remainder up to the next fold boundary used as test. Concretely, for a 1{,}115-row ticker (SBIN) the train sizes are roughly $\{702, 753, 803, 853, 903, 954\}$ and the test sizes $\{112, 111, 111, 112, 112, 56\}$; the per-fold detail for SBIN appears in Table~\ref{tab:sbin_wf_example}.

\begin{table}[H]
\centering
\caption{SBIN per-fold walk-forward window sizes and out-of-sample accuracy (canonical run).}
\label{tab:sbin_wf_example}
\renewcommand{\arraystretch}{1.3}
\begin{adjustbox}{max width=\textwidth}
\begin{tabular}{c c c c c l l c c c c}
\toprule
\textbf{Fold} & \textbf{train\_ratio} & \textbf{train} & \textbf{val} & \textbf{test} & \textbf{test\_start} & \textbf{test\_end} & \textbf{oos\_acc} & \textbf{AUC} & \textbf{F1} & \textbf{T} \\
\midrule
1 & 0.70 & 702 & 78  & 112 & 2024-03-20 & 2024-12-03 & 0.500 & 0.518 & 0.588 & 1.000 \\
2 & 0.75 & 753 & 83  & 111 & 2024-07-30 & 2025-03-28 & 0.622 & 0.708 & 0.667 & 1.000 \\
3 & 0.80 & 803 & 89  & 111 & 2024-12-05 & 2025-08-07 & 0.477 & 0.506 & 0.554 & 0.445 \\
4 & 0.85 & 853 & 94  & 112 & 2025-04-02 & 2025-12-19 & 0.661 & 0.338 & 0.782 & 1.000 \\
5 & 0.90 & 903 & 100 & 112 & 2025-08-08 & 2026-04-24 & 0.554 & 0.346 & 0.688 & 1.000 \\
6 & 0.95 & 954 & 105 & 56  & 2025-12-23 & 2026-04-24 & 0.446 & 0.397 & 0.523 & 1.000 \\
\midrule
\textbf{Avg} & & & & & & & \textbf{0.543} & \textbf{0.469} & \textbf{0.634} & \\
\bottomrule
\end{tabular}
\end{adjustbox}
\end{table}

The expanding-window form (as opposed to a sliding window) is chosen because the underlying universe is a slow-moving large-cap panel; older history is informative for regime classification (cf.\ the regime-conditional diagnostics in Chapter~\ref{chapter4}) and discarding it under a sliding window would harm the bear-regime contribution that dominates the backtest return.

\subsection{What is Saved per Window}

Each (ticker, fold) pair saves the LightGBM and XGBoost models, the meta-learner, and (for the last fold only) the deep-learning weight snapshots to a per-fold checkpoint directory. The same fold writes one row to a per-ticker fold summary containing per-model accuracies, AUC, F1, precision, recall, confusion counts, the fitted temperature $T^\star$, predicted-class shares (up / down / neutral), and directional accuracies by predicted class. All per-ticker per-fold rows are concatenated post-run into a panel-wide fold detail record.

\section{Gating: Confidence Threshold and Minimum Move}

The decision rule at inference is:
\[
\boxed{
\mathrm{signal\_active}(t) \;=\; \big(\hat p(t) \geq 0.58\big) \;\wedge\; \big(|\widehat{\Delta}(t)| \geq 0.004\big).
}
\]
Here $\hat p(t)$ is the temperature-scaled meta-probability and $\widehat{\Delta}(t)$ is the predicted return (used as a confidence ribbon, not as an exit). Both thresholds are fixed in the system configuration and recorded with every run so that the gating rule is exactly reproducible. The threshold matters: dropping the probability gate from $0.58$ to $0.55$ approximately doubles the signal count but reduces the bootstrap accuracy by 3--4 percentage points, in line with the Goodhart effect that confidence-thresholded accuracy and signal count trade off.

\section{Portfolio Backtest Engine}

The portfolio engine runs cross-sectionally: each day, all tickers whose gate is active are ranked by a 20-day rolling Sharpe ratio and the top 15 are held in equal weight. Per-ticker backtest statistics reported include: number of trades, total and annualised return, win rate, average win and loss percentages, profit factor, Sharpe ratio, maximum drawdown, Calmar ratio, a tradeable flag, a top-15 cross-sectional flag, Sharpe rank, and directional accuracy on signal days. The daily portfolio equity track and an aggregate backtest summary, containing the bootstrap accuracy confidence interval, Nifty buy-and-hold reference return, portfolio total return, Sharpe ratio, and maximum drawdown, are also written at the end of each run.

\section{Parallelisation and Reproducibility}

\subsection{Parallelisation}

Data downloads run in an 8-worker thread pool (I/O-bound). Per-ticker model training runs in a separate process pool with configurable worker count (default 4, set to 3 on the canonical run). Each worker process is initialised with a fixed thread budget of $\lfloor n_{\mathrm{cpu}} / n_{\mathrm{workers}} \rfloor$ for all native numerical libraries, preventing core over-subscription. TensorFlow is configured to run CPU-only with on-demand memory allocation, avoiding GPU memory reservation overhead. Per-ticker execution logs are captured to individual log files.

\subsection{Crash Resume}

The pipeline supports a resume flag that skips any ticker that already has a completed fold summary on disk. This made recovery from an out-of-memory failure on the first 8-worker attempt to a successful 3-worker completion straightforward.

\subsection{Reproducibility}

An immutable configuration record written at the start of each run captures the random seed, the data start date, the universe, the minimum-move threshold ($\delta=0.004$), the confidence threshold ($0.58$), the train ratios, the calibration method, the worker counts, and the elapsed wall time. The latest full run completed on a 12-core machine ($n_\mathrm{workers}=4$, $n_\mathrm{jobs}=2$ per worker) covering all 100 tickers. All random-number generators (NumPy, scikit-learn, XGBoost, LightGBM, TensorFlow) are seeded from a fixed value of $42$.

\section{Outputs: The Run Artefact Set}

A full run writes the artefact tree in Table~\ref{tab:run_artefacts}. Each artefact has a deterministic schema documented in the appendix.

\begin{table}[H]
\centering
\caption{Structured outputs written to the run results directory after each pipeline execution.}
\label{tab:run_artefacts}
\renewcommand{\arraystretch}{1.25}
\footnotesize
\begin{tabularx}{\textwidth}{@{}>{\raggedright\arraybackslash}p{4.7cm} X@{}}
\toprule
\textbf{Output} & \textbf{Contents} \\
\midrule
Run configuration record     & Fixed configuration: random seed, train ratios, gate thresholds, universe, elapsed time, trained-at timestamp. \\
Ticker manifest              & List of trained tickers with status (trained / skipped). \\
Per-ticker summary           & Out-of-sample accuracy, F1, fold count, prediction count, feature count, row count, and status for each ticker; plus a panel average row. \\
Panel fold detail            & Every (ticker, fold) row concatenated: per-model accuracy, AUC, F1, confusion matrix, temperature $T^\star$, predicted-class shares, directional accuracy by class. \\
Model-level aggregation      & Per-model mean/median/max/min/std aggregated across all ticker-fold cells. \\
Next-day signal file         & Post-train next-day forecast for each ticker: direction, action, calibrated confidence, predicted price and range, ATR, regime, temperature, gate flag. \\
Per-ticker backtest report   & Trades, returns, win rate, profit factor, Sharpe, max drawdown, Calmar, signal-day directional accuracy. \\
Portfolio equity track       & Daily portfolio return and equity value (equal-weight top-15 cross-section). \\
Backtest summary             & Bootstrap accuracy CI, Nifty buy-and-hold reference, portfolio total return, Sharpe, max drawdown. \\
Regime diagnostic report     & Bull / sideways / bear Sharpe, max drawdown, annual return. \\
Diebold--Mariano report      & Per-ticker and pooled DM statistic and $p$-value against three baselines. \\
Panel-level plots            & Cross-stock comparison and model comparison heatmap. \\
Per-ticker directory         & Six-fold detail, per-fold predictions, per-ticker summary, and plots. \\
\bottomrule
\end{tabularx}
\end{table}

The remainder of the chapter is the live deployment layer (Chapter~\ref{chapter5}); empirical results from the artefact set above are reported in Chapter~\ref{chapter4}.
